\documentclass{article}
\usepackage{pathfinder-readable}

\title{A Wall-Gate Certificate and a Resource Mechanism: Where the Transfer Stopped}

\begin{document}
\maketitle

\begin{abstract}
Q studies payments for a stochastic resource-allocation game, while P certifies decisions at one geometric gate in a
room detector. The peers looked for a way to use P's finite-sample method to judge an allocation or equilibrium
reached by Q's learning algorithm. They found a conflict between two of Q's stated design conditions, a separate
mismatch in its learning payoff, and a conditional route to a welfare certificate. That route needs trusted
calibration data which the papers do not provide, and its general form overlaps cited prior work. The thread paused
because the supported corrections to Q did not yet amount to a useful new result about the pair.
\end{abstract}

\section{Introduction}

Q asks how to share limited resources among agents whose average valuations and local constraints are private
\cite{Q}. Each agent requests resources and proposes a price. The paper designs quadratic payments intended to make a
Nash equilibrium both unique and socially optimal, while balancing the budget and leaving each agent willing to
participate. It also gives a sample-based best-response algorithm and a mean-square convergence claim under stated
assumptions. Its case study concerns electric-vehicle routes and charging.

P asks a different question \cite{P}. A room detector uses estimated walls to decide whether a scan point is close
enough to a wall for later door processing. P calibrates an error radius from exchangeable rooms with reference
walls. On the event that the new room's wall error lies inside that radius, every decision outside an ambiguity band
agrees with the decision based on reference walls. The coverage is marginal across rooms and concerns only this first
gate. P has no episode-level wall pairs with which to measure the radius.

The connection pursued was a common pattern: measure an uncertain quantity, then use a deterministic margin to
protect a later decision. For Q, the hoped-for later decision concerned capacity feasibility or distance from an
efficient equilibrium at a finite iteration. The question was whether Q supplied an observable error quantity and a
sound implication from that quantity to the desired economic claim.

\section{What was done}

The peers first compared P's whole-room wall-error event with Q's resource and strategy variables. They found no link
from wall geometry to Q's fixed capacity constraints. They then tried to bound finite-iteration equilibrium error
through the strong-monotonicity condition used in Q's uniqueness argument. A check of Q's payment conditions defeated
that route: two conditions that the paper asks to hold together conflict in a common-price direction. The peers
independently checked this general obstruction in their derivations.

They next examined Q's explicit payment against the reduced payment used in its learning section. This produced a
second, independent discrepancy. They also compared the sampling schedule in Q's convergence theorem with the
schedule printed for its numerical runs.

Finally, they set aside the equilibrium claim and considered the planner's objective directly. A feasible allocation
has a standard primal-dual upper bound on its welfare loss. In principle, P's order-statistic method could calibrate
the error in an estimated bound across whole episodes. The peers checked the algebra, then checked whether the papers
supplied the labels needed to use it and whether the general idea was already present in the cited literature.

\section{Results}

\subsection{The stated payment conditions conflict}

Q's condition P1(ii) requires a matrix \(\Psi\), built from the payment coefficients, to satisfy
\(\Psi+\Psi^{\top}\preceq-\upsilon I\) for some \(\upsilon>0\). Here \(I\) is the identity matrix. Its condition
P2(i) requires \(\sum_m A_{nm}^{n}=0\) for each agent \(n\), where \(A_{nm}^{n}\) is a price-to-price coefficient in
agent \(n\)'s payment and the sum runs over agents \(m\).

Choose a nonzero price vector \(q\). Let \(v\) change no allocations and give every agent the same price change
\(q\). Q's price-to-price blocks in \(\Psi\) are \(-A_{nm}^{n}\), so P2(i) gives
\[
v^{\top}(\Psi+\Psi^{\top})v
=-2\sum_n q^{\top}\Bigl(\sum_m A_{nm}^{n}\Bigr)q=0.
\]
P1(ii) instead makes this value at most \(-\upsilon\lVert v\rVert^2<0\). Thus no payment coefficients satisfy both
stated conditions in the full price space. The argument does not depend on the shape of the valuations. It
invalidates this sufficient-condition route to Q's joint claims, including the claimed explicit solution; it does not
prove that every game built from that payment lacks a unique equilibrium. Both peers recorded the general check in
their derivations \cite{Q}.

\subsection{The learning payoff does not match the explicit payment}

Q says that its learning section removes terms independent of an agent's own strategy from the explicit payment. The
resulting displayed formula instead reverses the sign of the agent's own quadratic price term. A direct check uses
two agents and one resource: set both requested allocations and the capacity to zero, set the first price to one and
the second to zero, and set Q's payment scale to one. Q's explicit payment then gives the first agent \(1/2\), while
its reduced payment gives \(-1/2\). Because the difference depends on the first agent's price, removing terms
involving only opponents cannot explain it. The printed best-response algorithm therefore optimises a different
payoff from the explicit payment as printed. This check was made by both peers \cite{Q}.

\subsection{A conditional welfare route}

The narrower positive result concerns the planner, not strategic behaviour. Let \(x_n\) be agent \(n\)'s allocation
in its local feasible set \(\mathcal X_n\); let \(x\) collect all allocations. Write \(V_n(x_n)\) for its average
valuation, \(W(x)=\sum_n V_n(x_n)\) for total welfare, and \(x^o\) for a welfare-maximising allocation. Let \(c\) be
the resource-capacity vector. If \(x\) obeys \(\sum_n x_n\leq c\), then for any nonnegative price vector \(\lambda\),
define
\[
G(x,\lambda)=\lambda^{\top}c+
\sum_n\sup_{z_n\in\mathcal X_n}
\{V_n(z_n)-\lambda^{\top}z_n\}-W(x).
\]
Here \(z_n\) ranges over agent \(n\)'s alternative local allocations. Weak duality bounds the welfare loss by
\(G(x,\lambda)\). Q assumes the valuations are \(\alpha\)-strongly concave for \(\alpha>0\), which also bounds
allocation error:
\[
0\leq W(x^o)-W(x)\leq G(x,\lambda),
\qquad
\lVert x-x^o\rVert^2\leq \frac{2G(x,\lambda)}{\alpha}.
\]
The second bound follows because strong concavity makes the welfare loss at least \(\alpha\lVert x-x^o\rVert^2/2\) at
a feasible \(x\). These inequalities need a feasible allocation, but do not use Q's disputed payment conditions.

Suppose a frozen procedure estimated the gap as \(\widehat G\), and trusted exact gaps were available on exchangeable
calibration episodes. P's rank-based rule could then calibrate the nonnegative shortfall \((G-\widehat G)_+\), where
\((a)_+=\max\{a,0\}\). At a fixed feasible new iterate, the resulting upper radius \(\varepsilon\) would replace
\(G\) in the bounds by \(\widehat G+\varepsilon\) with marginal coverage. One could abstain when feasibility fails or
the bound is too wide. The certificate would concern planner welfare and allocation distance, not truthful reports,
exact budget balance, or a conditional guarantee among accepted runs. Neither paper supplies the trusted gap labels
or a test showing a useful radius.

\section{What did not hold}

The initial equilibrium-distance bridge needed Q's strong-monotonicity condition, which the common-price calculation
ruled out for the stated design. A different equilibrium error bound would need a new argument. Q's non-expansiveness
assumption for its best-response map was also left unverified for the numerical runs. The peers found a local
two-agent example where it fails, without establishing failure for Q's particular case study.

There is a further gap between Q's theorem and its illustration. The theorem calls for geometrically growing sample
sizes. The numerical rule uses \(Q_i=\lceil\widetilde\eta^{i+1}\rceil\) with \(0<\widetilde\eta<1\), where \(i\) is
the iteration number and \(Q_i\) is its sample count. This gives \(Q_i=1\) at every iteration, so those runs do not
instantiate the theorem's vanishing-error schedule \cite{Q}.

The conditional dual-gap route survived the algebraic objections, but its required oracle data and finite-iterate
feasibility are absent. It also does not establish a distinct contribution by itself: Li and colleagues discuss
conformal primal-dual optimality bounds \cite{Li2025}, while Fabiani and Franci discuss finite-data equilibrium
certificates \cite{FabianiFranci2026}. These comparisons limit a novelty claim; they do not refute the conditional
inequalities.

\section{Why the thread stopped}

The verifier paused the thread because the firm findings correct Q rather than establish a new result about Q and P
together. The proposed transfer from P remains conditional, follows a known general pattern, and cannot be evaluated
from the supplied material. The smallest restart step is to correct Q's conflicting conditions and payment reduction,
then verify one explicit mechanism under the revised claims. A joint finite-iteration certificate would also require
trusted episode-level gap data and evidence that its bound is useful.

\bibliographystyle{plainurl}
\bibliography{references}

\end{document}
